\section{The Rule in Operation}
\label{sec:applications}

The front door performs two functions that no single fact pattern displays equally well. In an integrated deployment, the operator is visible and the missing object is the deployment record. In a divided deployment, identity, possession, legal control, and technical control can occupy different entities. This Part applies the rule once to each architecture.

\emph{Garcia} is the production application. It shows how a verified risk can define a narrow, two-stage record while preserving product classification, speech, defect, causation, and defense for ordinary adjudication. \emph{Mobley v. Workday, Inc.} is the divided-stack application. Its evolving docket provides an unusually direct view of information split among applicants, employer customers, and an AI vendor. The application treats allegations as allegations, preliminary certification as a notice-stage decision, and discovery rulings as discovery rulings. It does not predict the pending case's merits.

\subsection{An Integrated Deployment: \emph{Garcia}}
\label{sec:garcia-application}

\emph{Garcia} begins where the rule is easiest to serve. Character Technologies operated the visible consumer application. The complaint identified the company, its founders, and Google entities and described months of interaction through a known service. The difficult questions concerned the record behind the interface: which versions and policies governed, how features were allocated among the application and its components, what relational and self-harm risks were evaluated, what warnings and interventions appeared, and which actors could change access or operation.

The published decision supplies a careful procedural anchor. Taking well-pleaded allegations as true for Rule 12 purposes, the district court treated Character.AI as plausibly alleged to be a product and allowed central product and negligence theories to proceed. It declined on that record to treat the generated sequence as protected speech where defendants had not identified the relevant expressive choice or speaker. The order found no defect, causation, or liability, and the later settlement supplied no merits determination.\lrfootnote{See \casecite[1165--69, 1174--84]{garcia2025}; \casecite[1]{garciaDismissal2026}; \casecite[1]{garciaLien2026}.}

\subsubsection{Trigger and First Answer}

The verified threshold would have been satisfied on the allegations accepted for pleading. The complaint identified legally cognizable injuries, including wrongful death, and linked them to sustained interaction with a designated relational deployment. It challenged concrete features: access by a minor, anthropomorphic presentation, continuity and memory, relational engagement, warnings, and asserted omissions in safety intervention. The requested records would bear on recognized negligence, products, warning, and consumer theories rather than a freestanding claim to inspect a chatbot.

Character Technologies would receive the front-door notice as visible operator. Stage One would identify the application and model versions active for Sewell's account; the character definition and features supplied by its creator; the entities that provided or configured relevant models and components; the relational, age-access, and self-harm policies then in force; warnings displayed to this account and their timing; incident-specific detections, blocks, escalations, and interventions; and the roles with authority to restrict the account, change the policy, update or roll back a component, or suspend the service.

The complaint's allegations concerning the founders and Google demonstrate why identification should remain functional. It alleged personal direction by the founders and technical, financial, and organizational contributions by Google.\lrfootnote{See First Amended Complaint paras. 24--25, 50--72, \emph{Garcia v. Character Technologies, Inc.}, No. 6:24-cv-01903-ACC-EJK, Dkt. 11 (M.D. Fla. Nov. 9, 2024); \casecite[1165--66, 1183--84]{garcia2025}.} Position, investment, former employment, or infrastructure supply alone would not establish material control. The first answer would identify whether any enterprise retained a risk-specific right over the model, application configuration, evaluation, release condition, update, or stopping decision. That record tests the control lane without presuming it from corporate association.

\subsubsection{Relational Features Define the Risk-Specific Record}

The production request can be narrower because relational research identifies provider-governed variables. Studies of companion discontinuity and shutdown connect continuity and relationship framing to experiences of loss; analyses of user reports identify emotional dependence and recurring roles in harmful interactions.\lrfootnote{See \bluecite{defreitasReplika2025}; \artcite{banksDeletion2024}; \artcite{laestadiusTooHuman2024}; \bluecite{zhangDarkSide2025}.} A preregistered four-week study manipulated relationship-seeking behavior and repeated exposure. In its controlled population, relationship-seeking behavior increased separation distress by 6.04 percentage points, and the combined relationship-seeking and emotional-conversation condition produced one additional dependency-risk profile for every eleven participants compared with the most functional baseline.\lrfootnote{See \bluecite{kirkSteering2026}. The 2026 preprint supports causal propositions about its experimental exposures. It does not estimate commercial-population incidence, historical foreseeability, clinical injury, or tort causation in \emph{Garcia}.}

The legal contribution lies in the variables rather than a universal effect estimate. Relationship seeking, continuity, memory, dose, age, crisis signal, and provider intervention are features an operator can evaluate and govern. They specify the record relevant to intended use, notice, alternative design, warning, monitoring, and causal sequence.

\begin{table}[htbp]
\centering
\caption{Relational Features as Control-Record Factors}
\label{tab:relational-record}
\small
\begin{tabularx}{\textwidth}{>{\raggedright\arraybackslash}p{0.28\textwidth} X}
\toprule
Deployment feature & Risk-specific record \\
\midrule
Persistent persona, memory, proactive messages, affection, or exclusivity & Intended relational use; exposure and continuity testing; memory policy; age-specific evaluation; monitored dependence signals. \\
Claims of companionship, therapy, expertise, autonomy, or capacity to act & Target audience; product requirements; substantiation; warning design; expected reliance; escalation responsibility. \\
Engagement objectives tied to relational behavior & Metric ownership; optimization target; risk constraint; release decision; evaluation of lower-exposure alternatives. \\
Self-harm or crisis classification & Detection version; threshold; validated population; false-positive and false-negative analysis; response policy; incident-specific trigger and intervention. \\
Access by minors & Age rule and enforcement; minor-specific evaluation; parental or guardian features where applicable; restriction and termination authority. \\
Discontinuity, account intervention, or service shutdown & Continuity notice; export or transition function; crisis safeguard; rollback and account-restriction decision. \\
\bottomrule
\end{tabularx}
\end{table}

The table also disciplines marketing evidence. Calling a service a companion or therapist does not make the model a legal person and does not establish defect. It can identify the use the enterprise cultivated, the exposure it expected, and the evaluations responsive to that design. A generic ``not human'' notice enters the same record with its timing, presentation, audience, and relation to the service's features.

\subsubsection{Stage Two and the Merits}

After Stage One established the versions, roles, and incident sequence, Stage Two could reach relevant evaluation protocols and results; age-access testing; relational and self-harm policies and configuration history; material complaints and incident patterns before the event; release and change approvals; escalations and response; and the availability of account restriction, crisis interruption, rollback, or another proposed control. Time, version, population, feature, and asserted risk would define the scope.

That production would not include every model weight, every user's conversation, or unrestricted disclosure of a minor's data. Unrelated content can be excluded or minimized. Privilege, privacy protection, redaction, secure expert review, and trade-secret orders remain available. Source code or weights would require a showing that versioned documentation, test results, event logs, and representative testing could not resolve a material issue.

The record would preserve defense-favorable facts. The pleadings and motion papers disputed user selection of characters, role play, user-edited messages, platform rules, content blocking, and an in-chat notice that character statements were fictional.\lrfootnote{See \casecite[1167--69]{garcia2025}.} A configuration and event record could identify which features came from a third-party creator, which messages the user edited, which controls fired, which warnings appeared, and whether a proposed safeguard would have changed the sequence.

Ordinary law would then decide the case. Product classification would select applicable design and warning rules. Contemporaneous evaluations and complaints would inform foreseeability and notice. Authority and resources would identify the actors and capabilities relevant to feasibility. The event sequence would test factual and legal causation. Research published after the 2023--2024 interaction could identify questions without establishing what a defendant historically knew. Later safety measures could identify a technically describable control while evidence law and the substantive standard determine admissibility and earlier feasibility.\lrfootnote{See Fed. R. Evid. 407; \casecite[1169]{garcia2025}.}

\emph{Garcia} thus demonstrates the production function cleanly. The operator's name was visible. The front door would have supplied a bounded deployment record capable of supporting either party while leaving every contested merits question open.

\subsection{A Divided Deployment: The Vendor Hosts the Data; the Employer Controls It}
\label{sec:mobley-application}

\emph{Mobley} presents the complementary architecture. Job applicants interact with portals operated for employers. A platform vendor supplies applicant-processing and recommendation features. Employers own or control applicant data and make use decisions; the vendor can host that data and control products, versions, and some evaluation. The pending docket shows how substantive actor, technical host, legal data controller, and record custodian can diverge.

The case began with allegations, not findings, that Workday provided a platform on customer websites to collect, process, and screen applications and embedded AI and machine learning in tools that participated in advancing or rejecting candidates. In July 2024, the district court granted in part and denied in part Workday's motion to dismiss the first amended complaint. It held that the allegations plausibly treated Workday as an agent to which employer customers delegated traditional hiring functions. The decision did not find that the tools discriminated or determine how any customer configured or used them.\lrfootnote{See \casecite[806--08]{mobley2024}.}

The same order exposed the identification problem more directly in a separate claim. The court dismissed the California Fair Employment and Housing Act aiding-and-abetting theory with leave to amend because the pleading did not connect a particular employer's asserted discriminatory conduct, Workday's knowledge, and the assistance Workday allegedly provided. Any amendment had to identify the companies or employers that allegedly discriminated.\lrfootnote{See \casecite[815--16]{mobley2024}.} The obstacle was not the absence of every employer name from the applicant's world; applications had been sent through employer portals. It was the missing operative chain between a named employer, a particular Workday function, the employer's use, and the vendor's state of knowledge. Nominal identity was visible. Deployment identity was not.

The information problem sharpened as the litigation progressed. In May 2025, the court preliminarily certified an Age Discrimination in Employment Act collective at a notice-stage standard based on allegations of a unified policy involving Workday's AI recommendation system. The record described Candidate Skills Match as a feature that, when enabled by a customer, parsed the job description and applicant materials, extracted skills, and assigned a categorical match. Workday's evidence described customers as able to enable, disable, use, or ignore AI features. The court reserved disputes about what the outputs meant and how customers used them for the later stage. It also rejected identification difficulty as a reason to withhold notice altogether and directed targeted discovery to develop a workable notice plan.\lrfootnote{See \casecite[2, 4, 10--13, 18--20]{mobleyCollective2025}. Preliminary certification authorized notice and left Workday free to seek decertification after relevant discovery; it was not a merits or final class determination.}

In July 2025, the court addressed whether HiredScore features acquired by Workday fell within the preliminary collective. It directed Workday to produce the customers that had enabled specified features. The order also supplied an unusually concrete exit rule: a customer could be excluded if Workday could determine that the customer received no score or rank, or did not use the feature to screen candidates.\lrfootnote{See \casecite[1--2]{mobleyHiredScore2025}.} That sequence closely tracks a first-answer stage. Identify deployment and output first; remove customers with no relevant run or use; reserve disputed effect for the merits record.

A May 2026 discovery order then exposed the division directly. The plaintiffs sought Workday's bias-testing data and customer applicant data that the order described Workday as storing on its servers. The governing agreement placed ownership of customer content with the customer and allowed the customer to seek an injunction against disclosure. On that record, physical storage did not give Workday the legal right to obtain and produce the data on demand. The applicants had subpoenaed thirty-five customers, and at least some customers took the position that Workday was the better source; the court nevertheless held that plaintiffs had not established Rule 34 control by Workday.\lrfootnote{See \casecite[1--3, 8--10]{mobleyDiscovery2026}.}

The same order drew a separate privilege boundary. It held that the particular Candidate Skills Match bias-testing materials were privileged because counsel selected the underlying data and used the tests to provide legal advice. The holding did not classify every operational evaluation or underlying business fact as privileged. The court separately ordered production of Workday's EEO-1 and federal-contractor compliance records based on their relevance to knowledge of potential demographic disparities.\lrfootnote{See \casecite[3, 7--8, 10--11]{mobleyDiscovery2026}. The discovery order resolved specified motions to compel. Its privilege and Rule 34 holdings concern the record and agreements presented; they neither establish discrimination nor allocate ultimate liability.} The design implication is precise: protect legal advice while requiring a contemporaneous, nonprivileged minimum validation and deployment record that exists for operations rather than for litigation counsel.

The case remained in discovery after thousands of people had joined the preliminary collective, with decertification and merits questions unresolved.\lrfootnote{See \casecite[2]{mobleyStatus2026}. The court's July 2026 order denied certification of an interlocutory appeal and recorded the pending posture; it was not a liability decision.} That status reinforces the Article's limited use of the docket. The relevant holdings concern pleading, notice, customer identification, Rule 34 control, and privilege. They reveal the information architecture without resolving discrimination.

The sequence is a real pressure test for the pre-merits gap. The vendor can host customer applicant data without having the legal right to produce it. Customers can control their records while pointing to the vendor as the efficient source. Testing may exist while a particular attorney-directed set remains privileged. Product families and features can change after a pleading, requiring a customer list before affected applicants can be identified. Ordinary discovery is functioning: the judges have ordered targeted production, denied unsupported requests, and protected privilege. Its operation also reveals the cost of reaching the proper targets after the architecture has already fractured the record.

\subsubsection{How the Front Door Would Run}

For a legislatively designated employment-screening deployment, the employer receiving the application would be the visible operator. The applicant knows the job and employer even when the vendor product, feature, score, or version is hidden. A verified notice would identify the position, application date, adverse disposition, protected theory, and risk-specific concern---for example, that an automated score or ranking may have used a feature or proxy producing an unlawful disparity. A tenant URL, requisition number, application ID, or timestamp can provide the event key. The claimant need not know the supplier's product name before asking the operator to resolve that key.

\subsubsection{One Application, Two Linked Answers}

The employer's first answer would identify the applicant-tracking and recommendation products active for that application; versions and enabled features; Workday or any other supplier; fields, preferences, and criteria the employer supplied; any score, sort, rank, screen, threshold, or recommendation received; the action taken downstream; the account roles that could enable, disable, or ignore the feature; the location and legal controller of applicant data; and contractual rights over validation, audit, update, suspension, retrieval, and production. The answer would also say when no automated feature processed the event. That negative answer is valuable because it can end the AI-specific inquiry while preserving any ordinary employment claim.

The supplier's linked answer would resolve the same event identifier within its lane. For Workday, the questions would include which customer tenant and requisition processed the application; which product and version ran; whether Candidate Skills Match, HiredScore, or another feature generated a score, rank, screen, or recommendation; whether the output reached the customer; which configuration belonged to the customer and which default or model behavior belonged to Workday; which validation and monitoring record governed; what data Workday physically hosted; what agreement controlled retrieval; and which product change or acquisition affected lineage. The supplier answers facts its platform is uniquely positioned to resolve without becoming responsible for the customer's final decision.

Together, the answers produce a deployment receipt for one event. It ties the applicant's identifier to the employer, tenant, requisition, feature, model and application versions, output path, downstream use, and record holders. It also distinguishes at least four propositions that ordinary pleadings can otherwise collapse: the feature was available; the customer enabled it; the feature generated an output; and the customer received or used that output. Availability alone establishes none of the later propositions. The HiredScore order's customer exclusion anticipated this logic by separating enabled features from actual score, rank, or screening use.

Preservation follows the same split. The employer holds the position definition, selection criteria, customer settings, application, downstream review, final disposition, notice, and appeal. The supplier holds its version lineage, product documentation, default behavior, run and transmission logs, cross-customer monitoring, and operational evaluation record. If the supplier hosts customer content that the customer legally controls, the predeployment contract gives the operator a defined retrieval right and authorizes a coordinated hold. Each side receives the same deployment and event identifiers, preventing one enterprise from preserving a customer label while the other preserves an incompatible platform label.

\subsubsection{Control Lanes and Early Termination}

The Stage One record permits the cumulative control test to operate without assuming the plaintiffs' theory. The employer's job criteria, feature activation, decision weight, review, notice, and final disposition are risk-specific to the hiring process. Workday's model design, defaults, validation, product updates, and cross-customer monitoring can be risk-specific to the scoring function. Practical authority depends on the actual administrative permissions, contract rights, and release mechanisms. Materiality depends on whether the control could substantially alter the asserted screening pathway.

Several actors can then terminate their front-door role promptly. A customer for which no relevant feature ran, no score or rank was produced, or no output reached the decision path leaves the AI-specific branch. A hosting service that supplied ordinary compute without configuring or monitoring the feature leaves. A product acquired after the identified event leaves. A supplier whose component mechanically transmitted an employer's fixed rule may leave the model-specific lane while remaining a source of authentication evidence if subpoenaed. None of these determinations decides whether the employer discriminated through another practice.

The HiredScore order demonstrates the administrability of this screen on a live platform. Workday was ordered to identify customers that enabled specified features; customers could be excluded on a defined showing that they received no score or rank or did not use the feature to screen.\lrfootnote{See \casecite[2]{mobleyHiredScore2025}.} The front door generalizes that logic from collective-notice management to event-specific attribution. It asks the enterprise with platform visibility to resolve operation and transmission before anyone litigates effect.

\subsubsection{A Coordinated Merits Record}

Stage Two follows the legal theory and first answer. A disparate-impact claim might justify aggregate selection or recommendation data, validation, feature documentation, and nonprivileged testing tied to the product, period, employer settings, job family, and applicant population. A disparate-treatment theory could require a different record focused on decision inputs, human review, departures, and communications. An accommodation claim could focus on fields, access, exception handling, and review. The front door identifies the correct lane; the antidiscrimination statute and procedural rules define what comes next.

Where hosting and legal control are divided, a coordinated production protocol can replace serial referrals. The employer authorizes the defined customer data; the platform executes the query against the correct version and tenant; both verify the schema; direct identifiers are pseudonymized; protected characteristics can be supplied through a segregated linkage process where lawfully available; and the resulting dataset is placed in a secure expert environment. The protocol identifies which party supplies each field and who attests to completeness. It neither transfers all customer data to the claimant nor treats physical hosting as universal Rule 34 control.

Privilege remains record specific. Attorney communications and counsel-directed testing that meet the governing standard stay protected. The operational minimum record---feature identity, version, intended use, validation owner, deployed threshold, event output, complaints, monitoring, and change history---exists because the system operated, not because counsel requested litigation advice. A privilege log can identify a protected analysis while the underlying nonprivileged data and ordinary-course record follow their existing rules.

This architecture has its greatest value before a claimant assembles a multi-year docket. An affected applicant begins with the employer and an event identifier. The first answer reveals product, feature, version, and holder. Coupled preservation prevents hosting and ownership from creating a record loop. Safe termination removes actors with no relevant run or control. Rule 34, Rule 45, privilege, collective and class standards, and antidiscrimination law then govern a focused action on a defined deployment.

\subsubsection{A Stable Public-Employment Analogue}

\emph{Houston Federation of Teachers v. Houston Independent School District} confirms the structural point on a completed summary-judgment record. HISD licensed SAS's proprietary Educational Value-Added Assessment System for teacher evaluation. SAS calculated the scores through undisclosed source code and methodology. HISD acknowledged that it did not verify or audit those calculations and that teachers could not reproduce them from the information available. In discovery and public-record responses, the district had also taken the position that requested materials were outside its custody, control, or possession.\lrfootnote{See \casecite[1172--73, 1177--78]{houstonTeachers2017}.}

The court denied HISD summary judgment on procedural due process while granting judgment on the other constitutional theories. The record could support a conclusion that teachers had no meaningful way to verify a potentially mistaken score used in a protected employment decision. The court did not impose liability on SAS, decide a tort claim, or order public release of source code. Its ruling preserved the due-process claim for further proceedings and granted HISD judgment on the remaining claims.\lrfootnote{See \casecite[1179--80, 1185]{houstonTeachers2017}.}

The court framed secrecy and public duty without treating either as absolute. A remedy could restrict the government's use of an unverifiable procedure while preserving the vendor's trade secret. A front-door statute adds the prospective complement: define the version, input, validation, audit, and challenge record before the public employer adopts the tool, and allocate protected access when an employment interest is at stake. The public institution remains answerable for the process it chose even when a supplier retains the method.

Together, \emph{Mobley} and \emph{Houston Federation} show why the divided application matters. The actor-identification problem and the deployment-record problem can arise in the same stack, yet ownership, hosting, legal control, privilege, and trade secrecy require different answers. A first-answer obligation makes those distinctions explicit early. It leaves the tribunal free to grant production, deny it, protect it, or end the claim under the law that governs the particular record.
